\documentclass[a4paper,11pt]{article}
\usepackage[T2A]{fontenc}
\usepackage[utf8]{inputenc}
\usepackage[russian]{babel}
\usepackage{amsmath,amssymb,mathtools}
\usepackage{geometry}
\geometry{top=19mm,bottom=20mm,left=23mm,right=23mm}
\usepackage{microtype}
\usepackage{xcolor}
\usepackage{booktabs,tabularx,array}
\usepackage{tikz}
\usepackage{hyperref}
\usepackage{fancyhdr}
\usepackage{enumitem}
\usepackage{ragged2e}

\definecolor{accent}{HTML}{008080}
\definecolor{darkaccent}{HTML}{004D4D}
\definecolor{softgray}{HTML}{667176}
\definecolor{ink}{HTML}{273338}
\definecolor{error}{HTML}{A13E35}
\hypersetup{colorlinks=true,linkcolor=darkaccent,urlcolor=darkaccent,pdfauthor={Лёвин Артём Александрович},pdftitle={Ревью диагностической работы по физике — Софья Лузянина}}
\setlength{\parindent}{0pt}
\setlength{\parskip}{0.34em}
\setlength{\headheight}{15pt}
\setlength{\tabcolsep}{5pt}
\renewcommand{\arraystretch}{1.32}
\emergencystretch=2em
\pagestyle{fancy}
\fancyhf{}
\fancyhead[L]{\small\textcolor{softgray}{Физика \;\textbullet\; диагностическая работа}}
\fancyhead[R]{\small\textcolor{softgray}{15.11.2026}}
\fancyfoot[C]{\footnotesize\textcolor{softgray}{Лёвин Артём Александрович эксклюзивно для Софьи Лузяниной}}
\fancyfoot[R]{\footnotesize\textcolor{softgray}{\thepage}}
\renewcommand{\headrulewidth}{0.3pt}
\renewcommand{\footrulewidth}{0pt}
\newcommand{\themelabel}[1]{\par\vspace{0.20em}\noindent{\bfseries\textcolor{darkaccent}{#1}}\par\vspace{-0.42em}}
\newcommand{\errorlabel}[1]{\par\vspace{0.20em}\noindent{\bfseries\textcolor{error}{#1}}\par\vspace{-0.42em}}
\newcommand{\tasktitle}[2]{\phantomsection\label{task:#1}{\Large\bfseries\textcolor{darkaccent}{Задание №\,#1. #2}}\par\medskip\noindent\textcolor{accent}{\rule{\linewidth}{0.7pt}}\par\medskip}
\newcommand{\checkline}{\par\vspace{0.45em}\noindent\textcolor{accent}{\rule{\linewidth}{0.4pt}}\par}
\begin{document}
\thispagestyle{empty}
\vspace*{25mm}
{\color{accent}\rule{\linewidth}{1pt}}\par\vspace{16mm}
{\large\textcolor{darkaccent}{ФИЗИКА\quad/\quad ВХОДНАЯ ДИАГНОСТИКА}}\par\vspace{17mm}
{\fontsize{25}{32}\selectfont\bfseries Ревью диагностической\\[2mm] работы по физике}\par\vspace{14mm}
{\Large Софья Лузянина}\par\vspace{16mm}
{\large\textcolor{softgray}{15 ноября 2026 г.}}\par
\vfill
{\large Лёвин Артём Александрович}\par\vspace{6mm}
{\small\textcolor{softgray}{Разбор заданий, требующих повторения, и индивидуальное закрепление}}\par\vspace{25mm}
\begin{center}
\begin{tikzpicture}[x=1cm,y=1cm]
\draw[accent!60,line width=0.9pt] (0,0) .. controls (2,0.55) and (4,-0.45) .. (6,0.1)
.. controls (8,0.65) and (10,-0.1) .. (12,0.25);
\end{tikzpicture}
\end{center}
\clearpage

{\Large\bfseries\textcolor{darkaccent}{Задания для разбора}}\par\smallskip
В итоговый разбор включены только задания с отсутствующим решением. В основном диагностическом блоке это №\,18. Задания №\,29 и №\,30 относятся к отдельному опережающему блоку.
\medskip
\small
\begin{tabularx}{\linewidth}{@{}p{10mm}p{27mm}X p{23mm}p{36mm}@{}}
\toprule
\textbf{№} & \textbf{Тема} & \textbf{Что требует восстановления} & \textbf{Ваш ответ} & \textbf{Верный ответ}\\
\midrule
\hyperref[task:18]{18} & Влажность воздуха & Не выбрана формула влажности & «Не помню» & $50\%$\\
\addlinespace
\hyperref[task:29]{29} & Тонкая линза & Не применена формула линзы & «Не знаю» & $60$ см\\
\addlinespace
\hyperref[task:30]{30} & Энергия фотона & Не использована формула энергии фотона & «Не знаю» & $3{,}96\cdot10^{-19}$ Дж\\
\bottomrule
\end{tabularx}
\normalsize
\vspace{12mm}
{\bfseries\textcolor{darkaccent}{Как работать с разбором}}\par
Сначала восстановите логику решения по шагам. После этого закройте объяснение и самостоятельно выполните одну задачу на закрепление в конце соответствующего разбора.
\vspace{5mm}
\textcolor{softgray}{Формулировки «не помню» и «не знаю» означают, что ход решения не был представлен. Поэтому в разборе указаны недостающие действия, без приписывания ошибочных вычислений.}
\clearpage

\thispagestyle{fancy}
\vspace*{0.33\textheight}
\begin{center}
{\Huge\bfseries\textcolor{darkaccent}{Разбор ошибок}}\\[10mm]
{\large\textcolor{softgray}{Три задания \;\textbullet\; три способа решения}}\\[12mm]
\textcolor{accent}{\rule{0.62\linewidth}{0.7pt}}
\end{center}
\clearpage

\tasktitle{18}{Относительная влажность воздуха}
\themelabel{Текст задания}
При некоторой температуре плотность насыщенного водяного пара равна $17{,}3\ \text{г/м}^3$, а фактическая плотность водяного пара в воздухе равна $8{,}65\ \text{г/м}^3$. Найдите относительную влажность воздуха в процентах.

\themelabel{Ваше решение и ответ}
На листе написано: «не помню!». Вычислений и численного ответа нет.

\errorlabel{Что осталось невыполненным}
Последний верный шаг и первый ошибочный шаг определить невозможно: решение не было начато. Недостающий первый шаг --- выбрать определение относительной влажности через плотность водяного пара.

\themelabel{Ключевая идея}
При одной и той же температуре относительная влажность показывает, какую долю от максимально возможной (насыщенной) плотности пара составляет фактическая плотность:
\[
\boxed{\varphi=\frac{\rho}{\rho_{\text{нас}}}\cdot100\%.}
\]
Здесь $\rho$ --- фактическая плотность водяного пара, $\rho_{\text{нас}}$ --- плотность насыщенного пара \emph{при той же температуре}. У этих величин должны совпадать единицы измерения.

\themelabel{Полное правильное решение}
\begin{enumerate}[leftmargin=*,itemsep=0.16em]
\item Выписываем из условия $\rho=8{,}65\ \text{г/м}^3$ и $\rho_{\text{нас}}=17{,}3\ \text{г/м}^3$.
\item Делим фактическую плотность на насыщенную. Единицы $\text{г/м}^3$ сокращаются:
\[
\frac{\rho}{\rho_{\text{нас}}}=\frac{8{,}65}{17{,}3}=0{,}5.
\]
\item Переводим долю в проценты:
\[
\varphi=0{,}5\cdot100\%=50\%.
\]
\end{enumerate}

\themelabel{Проверка}
Фактическая плотность вдвое меньше плотности насыщенного пара. Значит, воздух содержит половину максимально возможного при данной температуре количества пара в единице объёма. Полученные $50\%$ физически разумны; также выполняется $0\leqslant\varphi\leqslant100\%$.

\themelabel{Итог}
\(\boxed{\varphi=50\%.}\)
\checkline
\themelabel{Задача на закрепление}
При одной и той же температуре плотность насыщенного водяного пара составляет $25{,}0\ \text{г/м}^3$, а фактическая плотность пара --- $15{,}0\ \text{г/м}^3$. Найдите относительную влажность воздуха в процентах.
\clearpage

\tasktitle{29}{Изображение в собирающей линзе}
\themelabel{Текст задания}
Предмет расположен на расстоянии $30$ см от собирающей линзы с фокусным расстоянием $20$ см. На каком расстоянии от линзы получится действительное изображение предмета? Ответ дайте в сантиметрах.

\themelabel{Ваше решение и ответ}
На листе написано: «не знаю». Вычислений и численного ответа нет.

\errorlabel{Что осталось невыполненным}
Вычислительную ошибку указать нельзя, поскольку решение не было начато. Отсутствует первый необходимый шаг: запись формулы тонкой линзы для собирающей линзы и действительного изображения.

\themelabel{Ключевая идея}
Свяжем фокусное расстояние $F$, расстояние от предмета до линзы $a$ и расстояние от линзы до действительного изображения $b$:
\[
\boxed{\frac1F=\frac1a+\frac1b.}
\]
В этой задаче собирающая линза имеет $F>0$, а действительное изображение располагается по другую сторону линзы, поэтому $b>0$.

\themelabel{Полное правильное решение}
\begin{enumerate}[leftmargin=*,itemsep=0.16em]
\item По условию $F=20$ см, $a=30$ см. Проверяем: $a>F$, поэтому действительное изображение возможно.
\item Выделяем неизвестное слагаемое:
\[
\frac1b=\frac1F-\frac1a.
\]
\item Подставляем величины в одинаковых единицах:
\[
\frac1b=\frac1{20}-\frac1{30}
=\frac{3-2}{60}=\frac1{60}\ \text{см}^{-1}.
\]
\item Следовательно, $b=60$ см.
\end{enumerate}

\themelabel{Проверка}
Подставляем ответ в исходную формулу:
\[
\frac1{30}+\frac1{60}=\frac2{60}+\frac1{60}=\frac3{60}=\frac1{20}.
\]
Равенство выполняется. Поскольку предмет находится между $F=20$ см и $2F=40$ см, изображение должно находиться дальше $2F$. Число $60$ см согласуется с этим качественным свойством собирающей линзы.

\themelabel{Итог}
\(\boxed{b=60\ \text{см}.}\)
\checkline
\themelabel{Задача на закрепление}
Предмет расположен на расстоянии $18$ см от собирающей линзы с фокусным расстоянием $12$ см. Определите расстояние от линзы до действительного изображения предмета.
\clearpage

\tasktitle{30}{Энергия фотона}
\themelabel{Текст задания}
Частота фотона равна $6{,}0\cdot10^{14}$ Гц. Найдите его энергию при $h=6{,}6\cdot10^{-34}\ \text{Дж}\cdot\text{с}$. Ответ запишите в единицах $10^{-19}$ Дж.

\themelabel{Ваше решение и ответ}
На листе написано: «не знаю». Вычислений и численного ответа нет.

\errorlabel{Что осталось невыполненным}
Последнего верного и первого неверного преобразования нет: ход решения отсутствует. Нужно вспомнить связь энергии фотона с его частотой и корректно выполнить умножение чисел в стандартном виде.

\themelabel{Ключевая идея}
Энергия одного фотона прямо пропорциональна частоте излучения:
\[
\boxed{E=h\nu,}
\]
где $h$ --- постоянная Планка, $\nu$ --- частота в герцах ($1\ \text{Гц}=1\ \text{с}^{-1}$), $E$ --- энергия в джоулях.

\themelabel{Полное правильное решение}
\begin{enumerate}[leftmargin=*,itemsep=0.16em]
\item Записываем формулу $E=h\nu$.
\item Подставляем данные:
\[
E=(6{,}6\cdot10^{-34})\,(6{,}0\cdot10^{14})\ \text{Дж}.
\]
\item Перемножаем коэффициенты и отдельно складываем показатели степеней:
\[
6{,}6\cdot6{,}0=39{,}6,\qquad -34+14=-20.
\]
\item Приводим результат к стандартному виду:
\[
E=39{,}6\cdot10^{-20}\ \text{Дж}
=3{,}96\cdot10^{-19}\ \text{Дж}.
\]
\end{enumerate}

\themelabel{Проверка}
Размерность: $\text{Дж}\cdot\text{с}\times\text{с}^{-1}=\text{Дж}$. Условие просит записать ответ \emph{в единицах} $10^{-19}$ Дж, поэтому в строку ответа записывается числовой коэффициент $3{,}96$.

\themelabel{Итог}
\(\boxed{E=3{,}96\cdot10^{-19}\ \text{Дж};\quad \text{число для ответа: }3{,}96.}\)
\checkline
\themelabel{Задача на закрепление}
Частота фотона равна $4{,}0\cdot10^{14}$ Гц. Принимая $h=6{,}6\cdot10^{-34}\ \text{Дж}\cdot\text{с}$, найдите энергию фотона. Ответ запишите в единицах $10^{-19}$ Дж.

\end{document}
